% Problem record op_c37d38eda1436c12. This TeX file is derived from
% database/problems_json/op_c37d38eda1436c12.json by scripts/sync-tex.mjs; edit the JSON record
% and rerun the script rather than editing this file.
\section{Information causality for almost-quantum correlations}
\label{sec:op_c37d38eda1436c12}

\paragraph{Problem.}
Do all almost-quantum correlations satisfy information causality?
A bipartite conditional distribution $p(a,b\mid x,y)$ with finite input and
output sets is almost quantum if there are a Hilbert space $\mathcal H$,
a unit vector $|\psi\rangle$, and projective measurements
$\{P_a^x\}_a$, $\{Q_b^y\}_b$ on $\mathcal H$ such that
\begin{equation}
  P_a^x Q_b^y|\psi\rangle=Q_b^y P_a^x|\psi\rangle,
  \qquad p(a,b\mid x,y)=\langle\psi|P_a^x Q_b^y|\psi\rangle.
  \label{eq:c37d-aq}
\end{equation}
Each measurement consists of orthogonal projections summing to the identity.
The commutation condition in \eqref{eq:c37d-aq} is imposed only on the shared
state, not as an operator identity.

Alice receives $L$ independent uniform bits $Z_1,\ldots,Z_L$.
Bob receives an independent uniform index $J\in\{1,\ldots,L\}$.
They may use shared randomness and almost-quantum boxes independent of these
inputs.
Alice sends Bob an $m$-bit classical message, with no communication
from Bob to Alice, and Bob outputs a binary guess $G$ for $Z_J$.
Information causality requires
\begin{equation}
  \sum_{j=1}^{L} I(Z_j:G\mid J=j)\leq m,
  \label{eq:c37d-ic}
\end{equation}
where $I$ is Shannon mutual information in bits.
Must \eqref{eq:c37d-ic} hold for every finite protocol of this form and all
integers $L\geq1$ and $m\geq0$?

\subsection*{Status}

Solved

\subsection*{Source}

Navascu\'es, Guryanova, Hoban, and Ac\'in define almost-quantum correlations
and leave their compatibility with information causality unresolved
(Definition~1 and Sections~9--10 of the arXiv version)
\sourcecite{ref:c37d-ngha}{NGHA15}.
The information-causality task and inequality originate with
Paw\l{}owski et al.\ \sourcecite{ref:c37d-ic-original}{PPKSWZ09}.

\subsection*{Progress}

\begin{itemize}
  \item Navascu\'es et al.\ verify several consequences of information causality
for almost-quantum correlations and report supporting numerical evidence,
while explicitly leaving the full principle unproved
\sourcecite{ref:c37d-ngha}{NGHA15}.

  \item Das Bhowmik and Banik give a negative answer in preprint version~1,
Corollary~1 \sourcecite{ref:c37d-bb}{BB26}.
Their Lemma~1 supplies a finite-dimensional almost-quantum box with $N$
settings and two outcomes per party.
A single use of this box and one
classical bit yield, for $L=kN$ independent data bits, the information score
\begin{equation}
  \sum_{j=1}^{kN} I(Z_j:G\mid J=j)
  \geq \frac{k}{17}\left(1-(1-2^{-k})^N\right).
  \label{eq:c37d-violation}
\end{equation}
For $k=18$ and $N=757693$, the right-hand side of
\eqref{eq:c37d-violation} exceeds $1.00000018$, contradicting
\eqref{eq:c37d-ic} with $m=1$.
Taking $N=k2^k$ instead gives the unbounded lower bound
$k(1-e^{-k})/17$.
The family uses increasing numbers of settings and data
bits. The claim is not an unbounded violation by one fixed finite box.
\end{itemize}

\subsection*{References}

\begin{enumerate}[label={},leftmargin=4.8em,labelsep=0.5em]
  \footnotesize
  \item[\textup{[NGHA15]}]\label{ref:c37d-ngha}
  Miguel Navascu\'es, Yelena Guryanova, Matty J. Hoban, and Antonio Ac\'in,
``Almost quantum correlations,'' \emph{Nature Communications} \textbf{6},
6288 (2015).
\href{https://doi.org/10.1038/ncomms7288}{doi:10.1038/ncomms7288};
\href{https://arxiv.org/abs/1403.4621v1}{arXiv:1403.4621v1}.

  \item[\textup{[PPKSWZ09]}]\label{ref:c37d-ic-original}
  Marcin Paw\l{}owski, Tomasz Paterek, Dagomir Kaszlikowski, Valerio Scarani,
Andreas Winter, and Marek Żukowski,
``Information causality as a physical principle,'' \emph{Nature}
\textbf{461}, 1101--1104 (2009).
\href{https://doi.org/10.1038/nature08400}{doi:10.1038/nature08400};
\href{https://arxiv.org/abs/0905.2292v3}{arXiv:0905.2292v3}.

  \item[\textup{[BB26]}]\label{ref:c37d-bb}
  Anandamay Das Bhowmik and Manik Banik,
``Almost-Quantum Correlations Violate Information Causality,'' preprint (2026).
\href{https://arxiv.org/abs/2610.41444v1}{arXiv:2610.41444v1}.
\end{enumerate}

\subsection*{Comment}

The universal information-causality claim is false by the finite
counterexample in \sourcecite{ref:c37d-bb}{BB26}.
The Solved designation records this negative resolution and its preprint source.
The score includes every trial.
Bob always supplies a binary guess, so the
violation does not arise from conditioning on successful certificates.

\subsection*{Field}

Quantum Foundations; Quantum Communication

\subsection*{Topic}

Bell nonlocality; Quantum communication complexity

\subsection*{ID}

\texttt{op\_c37d38eda1436c12}
